% !TEX root = ../main.tex
\chapter{Library: saved shapes and window sets}\label{ch:library}
The Library is the project's catalogue of what it has found and kept. It has three sections,
chosen with the tabs at the top of every Library page:
\begin{description}
  \item[Motifs] every shape the project has catalogued, every place each one occurs, and the
    families those shapes fall into. Five pages: \apppage{library.recurrence},
    \apppage{library.atlas}, \apppage{library.family}, \apppage{library.grouping} and
    \apppage{library.import}.
  \item[Window sets] saved, reusable sets of fixed-length windows cut from the recordings, the
    raw material for training a model (\apppage{library.window-sets}).
  \item[Templates] every saved chain (\apppage{library.templates}).
\end{description}
The Library feeds the rest of the app. A family's exemplar is the seed Discovery searches
for; a family is the source Analyse interrogates; a window set is the source of a training
chain; a template is a chain you can run again.

\begin{plainterms}[Motifs, members and families]
A \term{motif}{motif} is a shape, such as a sharp drop followed by a slow recovery. Each place
that shape occurs, on some channel at some time, is a \emph{member}. A \term{family}{family}
is a group of motifs that look alike, the way a field guide groups birds by silhouette. The
families are not fixed: they are the answer to a question you ask of the catalogue, such as
``group the single motifs by shape distance, cut at 0.6''. That question, with its answer, is
a \term{grouping}{grouping}. Ask it differently and you get different families; both answers
are kept, and the grouping bar always says which one you are looking at.
\end{plainterms}

\begin{plainterms}[Exemplar and medoid]
Each family has two representatives. The \term{medoid}{medoid} is computed: it is the member
that is, on average, closest in shape to all the others, the most typical one. The
\term{exemplar}{exemplar} is the family's chosen example. If a person has picked one with
\ui{Make exemplar}, that member is the exemplar; otherwise it is the member nearest the
medoid, which is usually the medoid itself. The exemplar is what Discovery searches for, so it
is worth choosing a clean one by hand.
\end{plainterms}

Every Motifs page carries the same two bars. The \emph{grouping bar} says which grouping is on
screen: what is grouped (\ui{single motifs}, \ui{sequences} or \ui{spike trains}), the basis
and cut (click it to switch to another saved grouping), how many motifs it left out
(\ui{omitted · flagged}, which opens a drawer listing them and why), and \ui{Edit grouping}.
The \emph{view bar} under it hides members you probably do not want to look at, without
deleting anything: \ui{noise floor on} (the default) hides members shallower than their
dataset's noise floor (Settings › Datasets; 0.1\,mV where none is set); \ui{fall} keeps
only members whose fall lasts between two durations in seconds; \ui{pure only} keeps only
windows that hold a single fall. The bar always says how many members it is hiding, and
\ui{show sub-floor} brings them back.

\begin{watchout}[Controls on the grouping bar that do nothing yet]
\ui{adjudicated only}, the \ui{$\ge$ N members} popover's \ui{Apply} and \ui{+ filter} show
``not wired yet''. (The minimum-members number does hide small families on the Atlas, but no
regrouping happens.) \ui{Find by shape} and \ui{Export} on the section bar, and
\ui{Export entry} on the Atlas and Family pages, are not wired either.
\end{watchout}

\pagesection{library.recurrence}{Recurrence}{Where each family occurs, hour by hour}{live}
Recurrence is a table with one row per family and one column per channel, grouped by
recording. Each cell says how often that family occurs on that channel. You use it to see
whether a shape is everywhere, or belongs to one channel or one recording.

\pagefield{When you'd come here}
It is the first page of the Motifs section. Come here to choose which channels to look at, or
to ask whether a family recurs across channels and recordings.

\screenshot{pages/library/library.recurrence--default}{Families down the side, channels across
the top, three recordings to a page. A ``?'' cell is on a channel nobody has reviewed; the
rail on the right describes the grouping on screen.}

\pagefield{Doing things here}
\begin{enumerate}
  \item \textbf{Read a cell.} \ui{per hour} (the default) shows members per hour of
    recording, on a colour scale that spans several decades; \ui{count} shows the raw
    number. A ``?'' means nobody has reviewed that channel, so ``absent'' and ``never looked''
    are told apart. A blank cell was reviewed and the family was not found. A number in an
    amber dashed cell means members are recorded on a channel nobody has reviewed. A red
    cell holds suspected cross-channel artifacts, kept visible.
  \item \textbf{Count with or without artifacts.} \ui{every member} counts everything.
    \ui{excluding artifacts} drops the members a person marked as artifacts in Review; a
    machine flag alone drops nothing. \ui{propagation counted once} also counts a travelling
    event once, on the channel where it starts. Both need the family to have been classified
    across channels on \apppage{library.family}; the page says how many have been.
  \item \textbf{Open a cell.} Click it to see that cell's members in a slideshow, in place,
    each drawn from the recording with context either side.
  \item \textbf{Build a selection.} Tick channels or recordings. The rail totals the
    channels, channel-hours, members and the share reviewed, and warns when two selected
    channels share ground (browsing both would double-count). \ui{Browse} \emph{N}
    \ui{channels} opens \apppage{library.atlas}.
  \item Page through recordings with the arrows beside \ui{recordings 1--3 of} \emph{N}.
\end{enumerate}

\pagefield{Where it leads}
\apppage{library.atlas}; \apppage{library.family} from a family's name.

\begin{underthehood}[Recurrence]
Client: \code{webui/client/src/library/RecurrencePage.tsx}, \code{chrome.tsx} (section bar,
grouping bar, view bar, omitted drawer), \code{CrossChannel.tsx} (the counting modes),
\code{MemberSlideshow.tsx}. Read: \code{getRecurrence} in
\code{webui/client/src/api/library.ts} calls \code{GET /api/library/recurrence}
(\code{webui/server/library.py}). The three counting modes are the core's, from
\code{family_recurrence} in \code{Working/library/matching.py}, and the rule in force is
printed behind the info icon beside them.
\end{underthehood}

\pagesection{library.atlas}{Atlas}{Families grouped by shape}{live}
The Atlas shows every family of the grouping as a card: its exemplar and medoid drawn in real
millivolts, its member count and how much of it has been judged. It is the page for seeing
which kinds of shape the project has, side by side.

\pagefield{When you'd come here}
From Recurrence with a selection, or directly to browse families. Come here to compare
shapes, to find a family worth a closer look, or to send one to Discovery or Analyse.

\screenshot{pages/library/library.atlas--family-selected}{One family selected: its members open
in a slideshow above the grid, and the rail on the right gives its exemplar and medoid, its
amplitude spread and its numbers.}

\pagefield{Doing things here}
\begin{enumerate}
  \item \textbf{Read a card.} Each card draws the exemplar (dark) and the medoid (family
    colour), with only their own baseline removed: no card is normalised, because the size
    of an event is part of the evidence. Each card has its own y axis so every shape is
    legible; the bar at its right marks the family's peak on one scale shared by the whole
    page, so you can still compare sizes. Badges give the number of recordings, hand edits,
    suspected artifacts, members the view hides, and the method's noise test.
  \item \textbf{Choose the order.} \ui{sort:} by \ui{family id}, \ui{members},
    \ui{judged fraction}, \ui{duration} or \ui{hand edits}. The scope chips say which
    channels came from Recurrence; remove one to widen the scope.
  \item \textbf{Open a family.} Click a card to select it and open its members in a
    slideshow above the grid; double-click to go to \apppage{library.family}. The arrow keys
    move the selection.
  \item \textbf{Use the rail.} It shows the exemplar and medoid with the distance between
    them, a histogram of peak-to-peak amplitude, duration, mean distance to the medoid,
    signal-to-noise ratio, how many are judged and the cross-channel counts. Its buttons:
    \ui{Open all} \emph{N} \ui{members}, \ui{Seed search in Discovery} (with the exemplar as
    the seed), and \ui{Interrogate in Analyse}.
  \item Switch to \ui{sequences} in the grouping bar to see families of sequences (runs of
    events in order) instead of single motifs.
\end{enumerate}

\pagefield{Where it leads}
\apppage{library.family}; \apppage{discovery.seed}; \apppage{analyse.interrogation}.

\begin{watchout}[Use the Family page to send members to Review]
The rail's \ui{Send} \emph{N} \ui{unjudged to Review} only shows a message saying it is not
wired; nothing is queued. The same button on \apppage{library.family} does create a queue.
The legend calls the exemplar ``human seed'', but unless someone has used \ui{Make exemplar}
it is chosen by the computer (the member nearest the medoid).
\end{watchout}

\begin{underthehood}[Atlas]
Client: \code{webui/client/src/library/AtlasPage.tsx}. Reads: \code{getMotifFamilies} and
\code{getSequenceFamilies} in \code{webui/client/src/api/library.ts}, calling
\code{GET /api/library/families} and \code{/sequence-families}. The exemplar and medoid traces
are decimated millivolts read from the channel files by \code{webui/server/library.py}; the
exemplar is a hand-made one if there is one, else the member with the smallest stored
distance to the medoid.
\end{underthehood}

\pagesection{library.family}{Family}{Exemplar and medoid, every member, hand edits}{live}
The Family page shows one family in full: a summary of its exemplar and medoid, every member
drawn from its own recording, and a rail with the selected member's details. It is where you
correct the computer's grouping by hand.

\pagefield{When you'd come here}
From a card on the Atlas, a family name on Recurrence, or a link from Review. Come here to
check that a family's members really are the same shape, to remove the ones that are not, to
choose a better exemplar, or to send the unjudged members to Review.

\screenshot{pages/library/library.family--member-selected}{A family with its medoid selected.
The summary at the top overlays ten members on the medoid; each member card below is drawn
from its own recording, sorted by distance to the medoid.}

\pagefield{Doing things here}
\subsection{Look at the members}\label{sec:library-members}
\begin{enumerate}
  \item The summary gives the counts: members, recordings and channels, how many are
    judged, hand edits, duration, depth and whether the family has been classified across
    channels. ``Judged'' means a member has a Review verdict, or matches a judged event, or
    falls inside reviewed windows; the info icon gives the rule.
  \item \ui{sort} orders the member cards by \ui{distance to medoid}, \ui{time},
    \ui{amplitude} or \ui{unjudged first}. \ui{hand edits only} shows only members someone
    edited by hand.
  \item Click a card to show the member in the rail: its recording, channel, onset and
    duration, how it was found, its revisions, its edges, its verdict and its tags. Verdicts
    are read-only here; \ui{Open in Review} goes to give one.
\end{enumerate}

\subsection{Correct the family by hand}\label{sec:library-hand-edits}
Hand edits are written to the database and survive a reload and a regrouping, because they
are tied to the shape (its content hash), not to the grouping.
\begin{enumerate}
  \item Tick members (or \ui{select page}), then \ui{Remove from family}, \ui{Add tag}. A
    removed member moves to a strip of removed members, where \ui{Restore} undoes it.
  \item With one member in the rail: \ui{Make exemplar}, \ui{Remove from family}, or
    \ui{+ tag}.
  \item \ui{Assign class} is disabled: the classes are not settled yet, so nothing is
    written.
\end{enumerate}

\subsection{Send the family on}\label{sec:library-family-send}
\begin{itemize}
  \item \ui{Send} \emph{N} \ui{to Review as a queue} makes a Review queue of the members not
    yet judged; afterwards the button opens it. A family of unjudged members is a proposal,
    not a finding.
  \item \ui{Seed search in Discovery} searches the recordings for more of the exemplar.
  \item \ui{Classify across channels} starts a job that asks, for every member, whether the
    same event appears on other channels at the same moment (a suspected artifact), a little
    later (propagation), or not at all. Unlike the quick view on
    \apppage{explore.cross-channel}, this tests each pair against chance. Its result feeds
    the counting modes on \apppage{library.recurrence}, and \ui{Send} \emph{N}
    \ui{suspected artifacts to Review} queues the flagged members for a person to decide.
\end{itemize}

\screenshot{pages/zzz_w_cross_channel/library.family--fixup-w-classify-across-channels}{The
family after classification across channels: members paired by bin, with the rule that
produced each count.}

\pagefield{Where it leads}
\apppage{review.inspector} for the family's queue; \apppage{discovery.seed};
\apppage{explore.cross-channel} from an edge's \ui{Explore} link, on exactly the window the
edge was measured on.

\begin{watchout}[Edits change distances]
Making a new exemplar, or redrawing a member's span, makes the stored distances to it out of
date. The page says ``edges partially stale until recomputed''; nothing recomputes them
automatically. \ui{Redraw in Explore} opens a page that still runs on example data
(\apppage{explore.span-edit}).
\end{watchout}

\begin{underthehood}[Family]
Client: \code{webui/client/src/library/FamilyPage.tsx}, \code{Edges.tsx},
\code{CrossChannel.tsx}. Read: \code{GET /api/library/family/{id}}. Hand edits:
\code{POST} and \code{DELETE /api/library/hand-edits}, re-applied on every read by
\code{apply_to_assignment} in \code{Working/library/hand_edits.py}; an edit names its family by
the family's medoid content hash, so it survives a regrouping that renames the family
(\code{docs/LIBRARY_STORAGE.md}, §8.1). Queues: \code{POST
/api/library/family/{id}/unjudged-queue} and \code{/suspected-artifacts}. Classification:
\code{POST /api/library/family/{id}/classify-cross-channel}, a \code{cross_channel} job.
\end{underthehood}

\pagesection{library.grouping}{Edit grouping}{Unit, basis, what does not fit, hand edits}{live}
Edit grouping is where you ask the catalogue a new question: group these motifs by this
measure, with these settings. It opens as a window over the page you came from. You compute
a preview first, and only then save the new grouping or apply it.

\pagefield{When you'd come here}
From \ui{Edit grouping} on any Motifs page, when the families on screen are too coarse or too
fine, or when you want to group by something other than shape: amplitude, duration,
frequency, a tag, or where the motifs came from.

\screenshot{pages/library/library.grouping--shape-distance-default}{The Edit grouping window for
a shape-distance grouping: the four numbered steps, then the measured preview, then Save and
Apply.}

\pagefield{Doing things here}
\begin{enumerate}
  \item \textbf{1 What to group:} single motifs, sequences or spike trains.
  \item \textbf{2 Group by:} nine bases in three kinds. \emph{Distance} bases (shape
    distance, sequence similarity) cluster by how alike two shapes are; the default method
    is Ward linkage with a \emph{cut}, and other registered methods (k-Shape, HDBSCAN) may be
    offered. \emph{Feature bins} sort each motif into a bin of one measured number (dominant
    frequency, peak-to-peak amplitude, duration, polarity), with quantile, log-spaced or fixed
    bin edges. \emph{Labels} group by tag or by provenance (recording, run or spike train). A
    basis that does not apply to the chosen unit is disabled, with the reason.
  \item \textbf{3 Parameters:} the settings of the basis, drawn over the real distribution of
    the measure where there is one. For Ward, the \emph{cut} is the most important: a lower
    cut makes more, tighter families; a higher cut makes fewer, looser ones.
  \item \textbf{4 What does not fit:} motifs that fit no family are always flagged and listed,
    never deleted. You can also omit and flag groups smaller than a minimum size (10 by
    default).
  \item \ui{Compute preview} runs the grouping as a job without writing anything. The preview
    gives the number of groups, members, omitted motifs (by reason), the cost of recomputing,
    and what will happen to every hand edit.
  \item \ui{Save as a grouping} writes it as a new saved grouping; \ui{Apply grouping} saves
    it and opens it in the Atlas. The grouping you were looking at stays saved either way.
\end{enumerate}

\begin{plainterms}[Why save groupings rather than overwrite them]
Two cuts can give families that look alike and are not. Keeping every grouping, each with its
own id, means a figure or a result can always say exactly which families it used, and you
can switch back. Hand edits are not part of any grouping: they are re-applied on top of
whichever grouping you open.
\end{plainterms}

\pagefield{Where it leads}
\apppage{library.atlas} with the new grouping.

\begin{watchout}[Large regroupings]
When the recompute would be too heavy to run here, \ui{Apply grouping} becomes \ui{Create
SLURM script}, which says it is not wired yet. Spike-train groupings can be saved but not
applied, because no Atlas draws them yet.
\end{watchout}

\begin{underthehood}[Edit grouping]
Client: \code{webui/client/src/library/GroupingPage.tsx}. Reads: \code{GET
/api/library/grouping-editor} (the bases, the distribution of each measure, and the reasons a
basis does not apply). Preview: \code{POST /api/library/groupings/run}, a \code{regroup} job
that writes nothing. Save: \code{POST /api/library/groupings}. The engine is
\code{run_grouping} in \code{Working/library/grouping/engine.py}; the omission reasons and the
surrogate null that flags noise-like families are described in
\code{docs/LIBRARY_STORAGE.md}, §6.
\end{underthehood}

\pagesection{library.import}{Import}{An empty library, and the motif import}{live}
Import brings motifs that were extracted outside the app into the Library. It always starts
with a dry run: the importer reads the bundle and reports what it would do, without writing
anything. The same page is what a new installation sees, because an empty library opens here.

\pagefield{When you'd come here}
When the Library is empty, or when a new bundle of extracted motifs is ready. The empty
library offers two ways in: judge candidates in Review (pressing \ui{S} creates an
exemplar), or import motifs you have already extracted.

\screenshot{pages/library/library.import--dry-run-populated}{A dry run of a registered bundle:
counts, the importer's checks, what the import would create, and a sample.}

\pagefield{Doing things here}
\begin{enumerate}
  \item Choose a \ui{bundle}. The list comes from the registry of datasets, not from a scan
    of the disk: register a bundle in Settings › Datasets and it appears here. A held-out
    bundle is listed but cannot be chosen.
  \item Read the dry run: the counts of motifs, spike trains, recordings and channels; the
    \emph{Checks}, each passed, warned or failed (click one for details); \emph{What the
    import creates}; and a \emph{Sample}.
  \item \ui{Import} \emph{N} \ui{motifs} starts the import as a job. Its progress shows in the
    footer and \ui{Cancel import} stops it. The import writes everything at the end, so a
    job that did not finish wrote nothing.
  \item When it is done, \ui{Open the library}.
\end{enumerate}
Importing is safe to run twice: members are keyed on their recording and span, so a second
run adds no duplicates. Imported entries arrive unjudged, and no grouping is computed by the
import.

\pagefield{Where it leads}
\apppage{library.grouping} to group the new motifs, then \apppage{library.recurrence}.

\begin{watchout}[Two parts of this page are not what they seem]
The sample thumbnails are not signal: they are drawings of a generic drop at each sample's
amplitude, and the caption says so. The tick box \ui{then send} \emph{N} \ui{to Review as a
queue} creates no queue (it shows a ``not wired yet'' message), yet the finished card says
``queue created for Review''.
\end{watchout}

\begin{underthehood}[Import]
Client: \code{webui/client/src/library/ImportPage.tsx}, \code{EmptyLibrary.tsx}. Reads:
\code{GET /api/library/import/bundles}; \code{POST /api/library/import/dry-run}, which runs the
real importer with \code{dry_run=True}. The import is \code{POST /api/library/import}, a job
followed through \code{GET /api/jobs/{id}}. Importers live in
\code{Working/library/importers/}; identity and duplicates are \code{docs/LIBRARY_STORAGE.md},
§2.
\end{underthehood}

\pagesection{library.window-sets}{Window sets}{Saved and reusable, with their split and train-safety}{live}
A \term{window-set}{window set} is a saved list of fixed-length \term{window}{windows} cut from
the recordings, for example every one-minute window of eight channels. Models are trained and
tested on window sets. This page lists every saved set, says how it was made and whether it is
safe to train on, and hands a set to Analyse.

\begin{plainterms}[The split, and why ``train-safe'' matters]
A model learns from some windows (the \emph{train} part) and is marked on others (the
\emph{test} part). If a test window sits right next to a training window, the two overlap in
time and share most of their signal, and the model is in effect marked on something it has
already seen: like setting an exam with questions copied from the homework. The app therefore
splits by time, in blocks (a \term{blocked-split}{blocked split}), and leaves a gap of at least
one window length between a training block and a test block. A set is \emph{train-safe} when it
has such a split with such gaps on every boundary, checked when the set was saved.
\end{plainterms}

\pagefield{When you'd come here}
Before training a model, to choose or build its windows; or to check how an existing set was
made and whether it is safe.

\screenshot{pages/zzzzz_af_window_pools/library.window-sets--fixup-af-the-new-sets-are-listed-with-scale-and-counts}{The
list of window sets with their scale, counts and check, and the selected set's details on the
right.}

\pagefield{Doing things here}
\subsection{Read a set}\label{sec:library-read-set}
\begin{enumerate}
  \item The table gives each set's name, source, \ui{window · stride · gap}, number of
    windows, split, share labelled, and its \ui{check}. Filter by check (\ui{train-safe},
    \ui{not train-safe}), recording or labelled share.
  \item Select a set. The rail shows who made it and its recipe hash, its check with the
    reason when it is not train-safe, a strip per channel showing which hours are train, test
    and gap, the spacing checks, and how many windows carry each label. \ui{coverage}
    switches between the labels \ui{now} and the labels \ui{at save}.
  \item A check other than train-safe says why: the gap is shorter than a window, there is no
    split, the set is a test sample (never trained on), the sampling rate was inferred and
    needs confirming in Settings › Datasets, or the roles are laid later, over a pool.
\end{enumerate}

\subsection{Build a set}\label{sec:library-new-set}
\ui{New window set} cuts unlabelled windows from one recording: choose the channels, one or
more scales (1, 10 and 30 minutes are ticked by default; each scale becomes its own set), and
whether to leave out spans marked as artifacts (they are counted). The windows do not overlap.
Building runs as a background job and the new sets appear in the list when it ends. These sets
carry no train and test roles of their own: the split is laid when sets are combined in the
\emph{Window pool} block in Analyse. \ui{New from Analyse} opens that chain instead.

\subsection{Use a set}\label{sec:library-use-set}
\ui{Use as source in Analyse} opens \apppage{analyse.chain} with the Window pool block as its
source and this set ticked. \ui{Re-split in Analyse} does the same for a set that is not
train-safe. \ui{Rename} renames the set in the database.

\pagefield{Where it leads}
\apppage{analyse.chain}; \apppage{models.launch} trains on labelled sets.

\begin{watchout}[What is not recorded]
\ui{used by} is empty for every set: nothing records which run or model used a window set, so
deleting cannot check that a set is unused. Deleting a set only hides it in this browser
session. \ui{Import} and \ui{Export} are not wired; the import window checks only the path
you type and says it has no dry run.
\end{watchout}

\begin{underthehood}[Window sets]
Client: \code{webui/client/src/library/WindowSetsPage.tsx}, \code{NewWindowSet.tsx}. Read:
\code{GET /api/library/windowsets} over the \code{window_sets} table
(\code{webui/server/library.py}). Rename: \code{PATCH /api/windowsets/{id}/name}
(\code{webui/server/shape_routes.py}). Building unlabelled sets: \code{POST
/api/windowsets/unlabelled} (\code{webui/server/training_routes.py}). Pools and the train/test
fence: \code{Working/training/pool.py}.
\end{underthehood}

\pagesection{library.templates}{Templates}{Every saved chain, its versions and scores}{live}
A \term{template}{template} is a saved \term{chain}{chain}: its blocks in order, with their
settings locked, so the same analysis can be run again on other data. This page lists every
template in the database and sends one where it can run.

\pagefield{When you'd come here}
To find a chain someone saved, see its stages and settings, and apply it to new recordings
in Discovery, open it in Analyse, or launch a training template in Models.

\screenshot{pages/library/library.templates--template-selected}{The template list with each
chain's stages as small glyphs; the rail shows the selected template's stages with their
locked settings.}

\pagefield{Doing things here}
\begin{enumerate}
  \item Filter by kind (\ui{all}, \ui{detection}, \ui{seed search}, \ui{training},
    \ui{interrogation}) or with \ui{contains a Model stage}; \ui{sort} by runs, name or
    version.
  \item Select a template. The rail lists its stages with their settings, its version, and its
    scores.
  \item \ui{Apply in Discovery} runs a detection or seed-search template across recordings
    (\apppage{discovery.runs}); \ui{Open in Analyse} opens it as a chain; a training template
    has \ui{Launch in Models} instead.
\end{enumerate}

\pagefield{Where it leads}
\apppage{discovery.runs}; \apppage{analyse.chain}; \apppage{analyse.interrogation};
\apppage{models.launch}.

\begin{watchout}[Versions and scores are not recorded]
Each template carries a version number, but nothing stores what changed between versions, so
the rail says ``Version history is not recorded'' and there is no diff. No run stores a
precision, recall or comparison with a null, so every template reads ``No scores are
recorded'', even one that has been run. The run count is found by searching saved run
configurations for the template's name, not by a direct link. \ui{Duplicate}, \ui{Import JSON}
and \ui{Archive} change this browser session only; \ui{Export JSON} shows the template's stages
and settings as JSON in a window; it is a description, not a recipe you can run.
\end{watchout}

\begin{underthehood}[Templates]
Client: \code{webui/client/src/library/TemplatesPage.tsx}. Read: \code{GET
/api/library/templates} (\code{get_templates} in \code{webui/server/library.py}), which maps
the one \code{templates} table; the run count is a \code{LIKE} match of the name inside
\code{configs.config_json}. Template writes for Analyse go through \code{/api/templates}
(\code{webui/server/analyse_routes.py}).
\end{underthehood}
